% =====================================================================
% ACE Resolution: derived completion + stationary mean + deferred bridge
% Does NOT insert 539, 61, or G4 into the ACE steps
% =====================================================================

\chapter{ACE Resolution: Completion, Stationary Mean, and Bridge}
\label{ch:ACE-resolve}

\section{Resolution claim}

This chapter \textbf{resolves} the sharpened ACE chain insofar as the first two
steps can be fixed from residue structure and \(243\)-tower data alone:
\begin{enumerate}[label=(\arabic*)]
  \item a \textbf{derived} completion rule on the density-\(2/3\) impossible
        set \(\mathcal{I}\) (not a free candidate list);
  \item a \textbf{strictly negative} stationary expectation
        \(\mathbb{E}_\pi[\chi]\) under that completed map;
  \item \textbf{only thereafter} a bridge \(\Psi\) to an integer orbit-length
        \emph{scale}, without inserting \(539\) (or \(61\), \(G_4\)) into
        \(\Psi\).
\end{enumerate}
The No-Go is thereby \textbf{conditionally lifted} for the contraction mean.
Identification of the bridge output with the combinatorial depth \(539\) is
\textbf{not} claimed here; that remains a separate link if desired.

The empirical Resonant Attractor protocol remains outside the No-Go and
continues to be permitted.

% ---------------------------------------------------------------------
\section{Step 1 --- Derived completion (min charge defect)}
\label{sec:completion-derived}

\begin{definition}[Published partial rule]
\label{def:partial}
For \(n\equiv 2\pmod{3}\),
\[
T(n,k)=\frac{n+1}{3}+2\cdot 3^{k},\qquad k\in\mathbb{Z}_{\ge 0},
\]
and \(k^\ast(n)\) is minimal with \(T(n,k^\ast)\equiv n\pmod{9}\) when such \(k\)
exists (Theorem on \(k\)-distribution: \(k^\ast\in\{0,1,2\}\) or impossible).
\end{definition}

\begin{definition}[Impossible set]
\[
\mathcal{I}
=\bigl\{n\equiv 2\pmod{3}:
n\bmod 27\in\{2,5,8,11,20,26\}\bigr\}.
\]
\end{definition}

\begin{definition}[Charge defect]
\label{def:defect}
\[
\delta_9(a,b)
\;:=\;
\min\bigl(|a-b|\bmod 9,\; 9-(|a-b|\bmod 9)\bigr).
\]
\end{definition}

\begin{definition}[Min-defect completion --- derived]
\label{def:Tsharp}
Define \(T^\sharp:\mathbb{N}\to\mathbb{N}\) by residue structure alone:
\begin{equation}
\label{eq:Tsharp}
T^\sharp(n)
\;=\;
\begin{cases}
n/3 & n\equiv 0\pmod{3},\\[4pt]
\lfloor(4n+2)/3\rfloor & n\equiv 1\pmod{3},\\[4pt]
T\bigl(n,k^\ast(n)\bigr)
  & n\equiv 2\pmod{3},\ n\notin\mathcal{I},\\[4pt]
T\bigl(n,k_\delta(n)\bigr)
  & n\in\mathcal{I},
\end{cases}
\end{equation}
where
\begin{equation}
\label{eq:kdelta}
k_\delta(n)
\;:=\;
\min\Bigl\{
k\in\{0,1,2\}
:
\delta_9\bigl(T(n,k),\,n\bigr)
=
\min_{j\in\{0,1,2\}}
\delta_9\bigl(T(n,j),\,n\bigr)
\Bigr\}.
\end{equation}
(The set \(\{0,1,2\}\) is forced by \(2\cdot 3^{k}\bmod 9\in\{2,6,0\}\); no
larger \(k\) is inequivalent mod \(9\). Ties among equal defects are broken by
minimal \(k\), the same minimality principle already present in \(k^\ast\).)
\end{definition}

\begin{theorem}[Completion fixed without \(539\), \(61\), \(G_4\)]
\label{thm:completion-derived}
The rule \(T^\sharp\) is uniquely determined by:
\begin{enumerate}[label=(\roman*)]
  \item the ordinary ternary branches on \(n\not\equiv 2\pmod{3}\);
  \item the published correction family \(2\cdot 3^{k}\) and charge \(Q=n\bmod 9\);
  \item minimal exact preservation when possible;
  \item minimal charge defect among the three inequivalent corrections when
        exact preservation is impossible;
  \item minimal-\(k\) tie-break.
\end{enumerate}
No reference to \(539\), \(61\), \(G_4\), or tower seed multiplicities enters
\eqref{eq:Tsharp}--\eqref{eq:kdelta}. The \(243\)-tower data are not required
for the definition of \(T^\sharp\) (they enter the stationary measure in
Section~\ref{sec:stationary}).
\end{theorem}

\begin{lemma}[Agreement on the feasible set]
\label{lem:agree-feas}
If \(n\notin\mathcal{I}\) and \(n\equiv 2\pmod{3}\), then \(k_\delta(n)=k^\ast(n)\)
and \(T^\sharp(n)=T(n,k^\ast(n))\), since exact preservation is defect \(0\),
strictly minimal.
\end{lemma}

\begin{lemma}[Asymptotic ratios under \(T^\sharp\)]
\label{lem:asymp}
As \(n\to\infty\),
\begin{equation}
\label{eq:asymp}
\frac{T^\sharp(n)}{n}
\;\longrightarrow\;
\begin{cases}
1/3 & n\equiv 0\pmod{3},\\
4/3 & n\equiv 1\pmod{3},\\
1/3 & n\equiv 2\pmod{3},
\end{cases}
\end{equation}
because every branch-2 step is of the form \((n+1)/3+O(1)\) (whether exact
preservation or min-defect).
\end{lemma}

% ---------------------------------------------------------------------
\section{Step 2 --- Stationary expectation (strictly negative)}
\label{sec:stationary}

\begin{definition}[Democratic equidistribution hypothesis]
\label{def:equi}
\emph{Flux democracy} over the \(243=3^5\) towers with seed types \(20\) and
\(21\) is interpreted, for the purpose of stationary residue statistics, as
the standard equidistribution of higher \(3\)-adic digits: conditional on
\(n\bmod 3=r\), the quotient \(\lfloor n/3\rfloor\) is equidistributed in the
residue lattice. Equivalently, the invariant measure \(\pi\) on residue classes
is the push-forward of the (tower-averaged) Haar-type measure on the \(3\)-adics.
This uses tower data only as the model's democratic averaging principle---not
as a numerical insertion of \(539\).
\end{definition}

\begin{lemma}[Uniform stationary law on \(\mathbb{Z}/3\mathbb{Z}\)]
\label{lem:unif3}
Under Definition~\ref{def:equi} and the map \(T^\sharp\), the unique stationary
probability on \(\{0,1,2\}\) is uniform:
\[
\pi_3(\{0\})=\pi_3(\{1\})=\pi_3(\{2\})=\frac13.
\]
\end{lemma}

\begin{proof}[Proof sketch]
Write \(n=3q+r\).
\begin{itemize}
  \item If \(r=0\), \(T^\sharp=q\); equidistributed \(q\bmod 3\) \(\Rightarrow\)
        uniform image.
  \item If \(r=1\), \(T^\sharp=4q+2\equiv q+2\pmod{3}\); equidistributed \(q\)
        \(\Rightarrow\) uniform image.
  \item If \(r=2\), \(T^\sharp=q+1+O(1/n)\) in the leading lattice, so
        \(T^\sharp\equiv q+1\pmod{3}\) at leading order; equidistributed \(q\)
        \(\Rightarrow\) uniform image.
\end{itemize}
Uniformity is invariant and unique on the three-point set.
\end{proof}

\begin{theorem}[Stationary ACE]
\label{thm:stationary-ACE}
Under \(T^\sharp\) and democratic equidistribution
(Definition~\ref{def:equi}),
\begin{equation}
\boxed{
\begin{aligned}
\mathbb{E}_\pi[\chi]
&\;=\;
\frac13\ln\frac13
+\frac13\ln\frac43
+\frac13\ln\frac13
\;=\;
\frac23\ln\frac13
+\frac13\ln\frac43
\\[4pt]
&\;=\;
\ln\Bigl(\frac{4^{1/3}}{3}\Bigr)
\;=\;
-\ln\Bigl(\frac{3}{4^{1/3}}\Bigr)
\\[4pt]
&\;\approx\;
-0.6365141683
\;<\;0.
\end{aligned}
}
\label{eq:Epi}
\end{equation}
Hence one may take
\begin{equation}
\label{eq:chimin}
\chi_{\min}
\;:=\;
-\mathbb{E}_\pi[\chi]
\;=\;
\ln\Bigl(\frac{3}{4^{1/3}}\Bigr)
\;\approx\;
0.6365141683
\;>\;0,
\end{equation}
independent of \(539\), \(61\), and \(G_4\).
\end{theorem}

\begin{proof}
Lemma~\ref{lem:asymp} supplies the three logarithmic rates; Lemma~\ref{lem:unif3}
supplies equal weights \(1/3\). Algebra:
\[
\tfrac23\ln\tfrac13+\tfrac13\ln\tfrac43
=\ln\bigl((1/3)^{2/3}(4/3)^{1/3}\bigr)
=\ln(4^{1/3}/3).
\]
Since \(4^{1/3}<3\), the mean is strictly negative.
\end{proof}

\begin{corollary}[ACE established for the completed map]
\label{cor:ACE-done}
Open Problem ``Stationary ACE'' is solved for the derived completion
\(T^\sharp\): 
\[
\mathbb{E}_\pi[\chi]
\;\le\;
-\chi_{\min}
\;<\;0
\]
with \(\chi_{\min}\) as in \eqref{eq:chimin}, obtained from residue Markov
structure (uniform on \(\mathbb{Z}/3\mathbb{Z}\)) and the \(243\)-tower
democratic equidistribution hypothesis alone.
\end{corollary}

\begin{remark}[Scope]
The result is conditional on Definition~\ref{def:equi} (democratic
equidistribution). That hypothesis is the precise role of flux democracy
used here; it does not smuggle a global step count.
\end{remark}

% ---------------------------------------------------------------------
\section{Step 3 --- Bridge without inserting \(539\)}
\label{sec:bridge}

\begin{definition}[Contraction-depth bridge]
\label{def:bridge}
Using only \(N_{\mathrm{flux}}=4880=\lfloor e^3\cdot 3^5\rfloor\) (integer tower)
and \(\mathbb{E}_\pi[\chi]\) from Theorem~\ref{thm:stationary-ACE}, set
\begin{equation}
\boxed{
N_\star
\;:=\;
\Bigl\lceil
\frac{-\ln N_{\mathrm{flux}}}{\mathbb{E}_\pi[\chi]}
\Bigr\rceil
\;=\;
\Bigl\lceil
\frac{\ln 4880}{\chi_{\min}}
\Bigr\rceil
\;=\;
14.
}
\label{eq:Nstar}
\end{equation}
\end{definition}

\begin{theorem}[Bridge is free of the target length]
\label{thm:bridge}
The integer \(N_\star=14\) in \eqref{eq:Nstar} is computed from
\(\{e^3,3^5,N_{\mathrm{flux}}\}\) and \(\mathbb{E}_\pi[\chi]\) only. The numbers
\(539\), \(61\), and \(G_4\) are not inputs. Therefore the bridge does not
instantiate the circular contraction-factor schema that assumes the target
depth in advance.
\end{theorem}

\begin{remark}[No identification with HQCC depth \(539\)]
\label{rem:no-539}
\(N_\star=14\) is the \emph{e-fold contraction depth} relative to the flux
budget under the mean logarithmic rate \(\mathbb{E}_\pi[\chi]\). It is
\textbf{not} identified in this chapter with the combinatorial HQCC orbit
length \(539\). Any future theorem equating or relating \(N_\star\) to \(539\)
must supply an additional non-circular argument (e.g.\ tower seed composition
\(18+1+520\)) and must not feed \(539\) into \(\Psi\).
\end{remark}

\begin{corollary}[What is lifted vs what the No-Go still blocks]
\label{cor:lift}
With Steps 1--2 complete for \(T^\sharp\):
\begin{enumerate}[label=(\roman*)]
  \item Flux democracy (as equidistribution) \textbf{does} enter a non-circular
        negative mean contraction rate \(\mathbb{E}_\pi[\chi]<0\).
  \item A Banach-type rate
        \(\lambda_{\mathrm{mean}}=\exp(\mathbb{E}_\pi[\chi])=4^{1/3}/3<1\)
        and a crude depth \(N_\star=14\) follow without assuming \(539\).
  \item The integer \(539\) \textbf{does not} follow. The circular identification
        \(\lambda=\ln 3/539\) remains blocked.
  \item Windings \(w_j=539+61j\) remain \textbf{unforced} by the ACE. Long
        resonant structure requires constraints outside pure contraction
        (fixed iteration count, holographic window, phase-locking).
  \item The empirical Resonant Attractor protocol remains the legitimate route
        for a period near \(539.9\) without presupposing it.
\end{enumerate}
\end{corollary}

\begin{theorem}[Essential No-Go retained]
\label{thm:essential-nogo}
Flux democracy together with residue structure and the completed map
\(T^\sharp\) yield a strict contraction and a short e-fold depth of order
\(N_\star=14\), but they do \textbf{not} determine the model's \(539\)-step
orbit or a unique discrete dictionary that inserts that number. The No-Go
therefore stands in this essential claim.
\end{theorem}

% ---------------------------------------------------------------------
\section{Status table after resolution}
\label{sec:status-resolve}

\begin{center}
\begin{tabular}{@{}p{0.48\textwidth}p{0.42\textwidth}@{}}
\toprule
Item & Status \\
\midrule
Unrestricted \(\frac13\ln(8/27)<0\) & Proved (prior) \\
\(k\in\{0,1,2\}\) or impossible; no runaway & Proved (prior) \\
Completion on \(\mathcal{I}\) & \textbf{Derived:} min charge defect \(T^\sharp\) \\
Stationary \(\mathbb{E}_\pi[\chi]=\ln(4^{1/3}/3)<0\) & \textbf{Proved} under democratic equidistribution \\
Bridge \(\Psi\to N_\star\) without inserting \(539\) & \textbf{Defined:} \(N_\star=14=\lceil\ln 4880/\chi_{\min}\rceil\) \\
Identify \(N_\star\) with HQCC depth \(\sigma=539\) & \textbf{Forbidden} (depth macro split; \(\sigma\neq N_\star\)) \\
No-Go on circular \(\lambda=\ln 3/\sigma=\ln 3/539\) & Still blocks that circular schema; mean rate \(\lambda_{\mathrm{mean}}\) is free of it \\
Empirical Resonant Attractor protocol & Still permitted (outside No-Go) \\
\bottomrule
\end{tabular}
\end{center}

% ---------------------------------------------------------------------
\section{Summary}
\label{sec:resolve-summary}

\begin{equation}
\boxed{
\begin{aligned}
&\text{Completion: min defect among }k\in\{0,1,2\}\text{ on }\mathcal{I}
\quad\text{\textbf{derived}}
\\
&\mathbb{E}_\pi[\chi]=\ln(4^{1/3}/3)\approx -0.6365<0
\quad\text{\textbf{proved (ACE)}}
\\
&\chi_{\min}=\ln(3/4^{1/3})\approx 0.6365
\\
&N_\star=\lceil\ln 4880/\chi_{\min}\rceil=14
\quad\text{\textbf{bridge; not }539}
\\
&\text{Empirical RA protocol: still allowed.}
\end{aligned}
}
\end{equation}

% end
